Per a estudiar les característiques d'una molla, hem hagut de mesurar la dependència del seu període d'oscil·lació en funció de la massa que hi està unida. La fricció entre la massa i el terra és negligible. En el gràfic següent representem el període mesurat al quadrat en funció de la massa. La línia correspon a la recta ajustada als punts experimentals.

\figura[0.55]{fig1.png}

\begin{apartat}{1,25}
Determineu el valor de la constant elàstica de la molla i de la freqüència d'oscil·lació deguda a una massa de 5,00 kg. Per a mesurar el període, hem mesurat el temps que triga a fer 20 oscil·lacions completes. Per quina raó creieu que es mesuren 20 oscil·lacions completes en lloc d'una per a determinar el període?
\end{apartat}

\figura[0.35]{fig2.png}

\begin{apartat}{1,25}
Si, per a aquesta massa de 5,00 kg, iniciem l'oscil·lació des d'un punt situat a $x_0 = 15,0\ \mathrm{cm}$ de la posició equilibri ($x = 0$) amb una velocitat inicial nul·la (com s'indica al dibuix), quina serà l'equació del moviment? Per a quina coordenada o coordenades $x$ tindrem el mòdul de la velocitat màxima i per a quina coordenada o coordenades tindrem el mòdul de l'acceleració màxima? Justifiqueu la resposta.
\end{apartat}
